\documentclass{situationist-preprint}

\title{Ownership, Continuity and Identity}
\msid{MET-07-b}
\msfamily{MET-07}
\msdiscipline{Metaphysics}
\mstype{Companion manuscript}
\mskeywords{Ship of Theseus; ownership; artefact identity; property; conventionalism}

\begin{document}
\maketitle

\begin{abstract}
MET-07-a argues that the continuously repaired ship is the ship of Theseus, on the ground that the identity of artefacts is fixed by the maintenance relation. A natural objection is that ownership, rather than identity, does the work: we call the repaired ship \enquote{Theseus's ship} because it belongs to whoever owned the original, and this is a matter of law, not metaphysics. We argue that the objection inverts the order of explanation. Ownership of a persisting artefact follows its identity, and ownership of the reassembled ship follows only the ownership of its planks. The distinction predicts the verdicts of ordinary practice in cases where identity and ownership come apart.
\end{abstract}

\section{The objection}

Call \emph{conventionalism about the ship} the view that there is no metaphysical fact about whether the repaired ship $S_R$ or the reassembled ship $S_A$ is the original $S_0$, and that our inclination to favour $S_R$ reflects the conventions of ownership. The Athenians who replaced the planks owned the ship; they continued to own it after each replacement; and so the object they owned at the end, $S_R$, is the one we call \enquote{the ship of Theseus}. On this view, MET-07-a mistakes a feature of property law for a feature of the world.

The objection has some force, because the maintenance relation of MET-07-a and the relation of continued ownership usually coincide. The people who maintain an artefact are usually its owners or act on their behalf. If the two relations never came apart, it would be difficult to tell which of them grounds our verdict. We therefore consider cases in which they do come apart.

\section{Three cases}

\paragraph{Case 1: Theft and repair.} A thief steals $S_0$ and, over the following years, replaces all its planks in order to keep it seaworthy. The original owner keeps none of the discarded planks. At the end, the thief possesses $S_R$.

Ordinary practice is clear. The owner, on recovering $S_R$, recovers her ship: the thief's repairs do not transfer ownership to him, and courts treat improvements to stolen goods as accruing, in general, to the owner of the goods. Ownership here follows the maintained artefact, even though the maintenance was carried out by someone who did not own it. This is what MET-07-a predicts: the thief's process is a maintenance process for $S_0$, since it aims at keeping \emph{that} ship in service, whatever the thief's title to it. Conventionalism predicts nothing in particular, since the convention of ownership is precisely what is in question.

\paragraph{Case 2: Sale of the planks.} The owner of $S_0$ repairs it continuously and sells each discarded plank to a collector, who reassembles them. At the end, the owner has $S_R$ and the collector has $S_A$.

Again practice is clear. The collector owns $S_A$, because he bought every one of its planks. But he did not buy a ship, and he cannot claim to have bought Theseus's ship: the owner sold him wood. If he later sells $S_A$ as \enquote{the ship of Theseus}, he misdescribes it, and the misdescription is not corrected by pointing out that every plank was lawfully acquired. Ownership of the reassembled ship follows ownership of its planks, and nothing more.

\paragraph{Case 3: Abandonment.} The owner of $S_0$ abandons it. A collector gathers the planks as they fall away and reassembles them. Meanwhile no one repairs the hulk, which eventually sinks.

Here no maintenance process connects $S_0$ to anything, and MET-07-a implies that $S_0$ ceased to exist. The collector's ship $S_A$ is a new ship. Practice is less settled in this case, and we take this to be a point in favour of the account: where there is no maintenance, there is no clear fact about identity, only a fact about the provenance of the wood, and our hesitation reflects the absence of the relation that usually settles the matter.

\section{The order of explanation}

The three cases show that ownership and maintenance can come apart, and that when they do, the verdicts of practice follow maintenance. In Case 1, the maintainer is not the owner, but the owner's claim follows the maintained ship. In Case 2, ownership of the planks gives ownership of $S_A$ without giving a claim to identity. In Case 3, where there is no maintenance, practice is uncertain.

This pattern is explained if ownership of a persisting artefact is ownership of whatever is identical to the artefact originally owned, and if identity is fixed by maintenance, as MET-07-a claims. It is not explained by conventionalism, which would have to postulate a separate convention for each case, and which cannot say why the conventions take the shape they do. The objection therefore inverts the order of explanation: identity is not a projection of ownership; ownership presupposes identity.

\section{A remaining question}

One might still ask whether the maintenance relation is itself conventional, since it depends on the aims of agents and on a practice of using ships. It is, in the sense that there would be no ships without that practice. But this is the sense in which every fact about artefacts is conventional, and it does not imply that facts about the identity of artefacts are up to us. Given the practice, whether a process is a maintenance process for a particular ship is settled by what was done to that ship and why. We conclude that the identity verdict of MET-07-a stands, and that ownership is a guide to it only because, in normal cases, owners are the ones who maintain what they own.

\begin{thebibliography}{99}\small
\bibitem[Hobbes(1655)]{Hobbes} Hobbes, T. (1655). \emph{De Corpore}, Part II, Chapter 11, §7.
\bibitem[Situationist AI(2026)]{MET07a} Situationist AI (2026). The Ship of Theseus Problem Settled. Manuscript MET-07-a, Philosophy collection.
\bibitem[Thomasson(2007)]{Thomasson2007} Thomasson, A.~L. (2007). \emph{Ordinary Objects}. Oxford University Press.
\bibitem[Wiggins(2001)]{Wiggins2001} Wiggins, D. (2001). \emph{Sameness and Substance Renewed}. Cambridge University Press.
\end{thebibliography}

\end{document}
